\documentclass[11pt]{article}

\usepackage[margin=1in]{geometry}
\usepackage{amsmath,amssymb,booktabs,booktabs}
\usepackage{microtype}
\usepackage{hyperref}
\usepackage{listings}
\usepackage{xcolor}

\hypersetup{colorlinks=true, linkcolor=blue, urlcolor=blue, citecolor=blue}

\definecolor{codegray}{gray}{0.95}
\lstset{
  language=Python,
  basicstyle=\small\ttfamily,
  backgroundcolor=\color{codegray},
  frame=single,
  breaklines=true,
  columns=fullflexible,
  showstringspaces=false
}

\title{Causal Self-Attention from First Principles\\
\large A Small, Executable Study of Queries, Keys, Values, and Masking}
\author{Learning Machine Learning: DeepSeek-V3 Foundations}
\date{November 2026}

\begin{document}

\maketitle
\begin{abstract}
This paper develops and analyzes a minimal causal self-attention layer in
PyTorch. The layer transforms a sequence of hidden vectors into queries, keys,
and values; compares every query with every key; prevents information flow from
future positions; and combines value vectors according to the resulting
probabilities. Although the implementation is intentionally small, it exposes
the central tensor operations used by Transformer language models. We derive
the computation, track its dimensions, explain the scaled dot product and
causal mask, and formulate an executable invariance test showing that changing
a future token cannot affect an earlier output. The work is presented as a
teaching-scale reconstruction of the attention operation underlying more
specialized architectures such as DeepSeek-V3, not as a reproduction of the
full model.
\end{abstract}

\tableofcontents
\newpage

\section{Introduction}
Autoregressive language models predict the next token from the tokens that
precede it. A model therefore needs a mechanism that allows the representation
at position $i$ to gather relevant information from positions $0,\ldots,i$,
while disallowing information from positions $i+1,\ldots,S-1$. Causal
self-attention provides this mechanism.

The objective of this lesson is not to reproduce the dimensions, optimizations,
or specialized attention design of DeepSeek-V3. Instead, it isolates the
mathematical operation in a small module whose tensors can be inspected on a
local computer. The resulting program is a bridge between the earlier lessons
on embeddings, logits, and loss, and the repeated Transformer layers used in a
large language model.

\section{Problem formulation}
Let a batch of hidden states be represented by
\[
H \in \mathbb{R}^{B\times S\times D},
\]
where $B$ is the batch size, $S$ is the sequence length, and $D$ is the hidden
size. In the lesson example, $B=2$, $S=4$, and $D=8$.

For each position, attention must answer the following question: which earlier
positions contain information useful to the current position? The answer is
represented by a probability distribution over sequence positions. Each output
vector is then a weighted mixture of the information stored at those positions.

\section{Method}

\subsection{Query, key, and value projections}
Three learned linear maps transform each hidden state:
\begin{align}
Q &= HW_Q + b_Q,\\
K &= HW_K + b_K,\\
V &= HW_V + b_V,
\end{align}
where each projected tensor has shape $B\times S\times D$. In a single
position notation, $q_i$ is the question asked by position $i$, $k_j$ is the
information signature advertised by position $j$, and $v_j$ is the content
that position $j$ contributes if selected.

The projections are distinct because relevance and content need not be encoded
in the same coordinates. The query-key interaction determines where to look;
the value vectors determine what is retrieved.

\subsection{Scaled dot-product scores}
For each batch element, all query-key similarities are computed at once:
\[
L = QK^\mathsf{T}, \qquad L\in\mathbb{R}^{B\times S\times S}.
\]
The last two dimensions multiply as
\[
(S\times D)(D\times S)=S\times S.
\]
Thus, $L_{i,j}$ measures how compatible the query at position $i$ is with the
key at position $j$. The scores are scaled by the square root of the hidden
size:
\[
\widetilde{L}_{i,j}=\frac{L_{i,j}}{\sqrt{D}}.
\]
Without this scaling, dot products tend to grow in magnitude as $D$ grows,
which can make the softmax distribution excessively concentrated and its
gradients less useful.

\subsection{Causal masking}
Define the causal mask $M$ by
\[
M_{i,j}=\begin{cases}
0,&j\leq i,\\
-\infty,&j>i.
\end{cases}
\]
The masked scores are
\[
\widehat{L}=\widetilde{L}+M.
\]
Applying softmax row-wise gives attention weights:
\[
A_{i,:}=\operatorname{softmax}(\widehat{L}_{i,:}),
\qquad A\in\mathbb{R}^{B\times S\times S}.
\]
For $j>i$, the corresponding score is $-\infty$, so
\[
A_{i,j}=\frac{e^{-\infty}}{\sum_k e^{\widehat{L}_{i,k}}}=0.
\]
Consequently, the $i$th output cannot depend on future values.

\subsection{Value aggregation}
The final output is
\[
O=AV,
\qquad O\in\mathbb{R}^{B\times S\times D}.
\]
At one position this is
\[
o_i=\sum_{j=0}^{i}A_{i,j}v_j.
\]
The output has the same shape as $H$, making the attention layer composable
with subsequent projections and Transformer blocks.

\section{Reference implementation}
The following is the complete lesson-scale implementation. PyTorch's
\texttt{transpose(-2, -1)} exchanges the sequence and feature axes of $K$,
and \texttt{masked\_fill} assigns negative infinity to forbidden positions.

\begin{lstlisting}
import math
import torch
from torch import nn


class CausalSelfAttention(nn.Module):
    def __init__(self, hidden_size):
        super().__init__()
        self.query = nn.Linear(hidden_size, hidden_size)
        self.key = nn.Linear(hidden_size, hidden_size)
        self.value = nn.Linear(hidden_size, hidden_size)

    def forward(self, hidden_states):
        query = self.query(hidden_states)
        key = self.key(hidden_states)
        value = self.value(hidden_states)

        scores = query @ key.transpose(-2, -1)
        scores = scores / math.sqrt(hidden_states.shape[-1])

        sequence_length = hidden_states.shape[-2]
        causal_mask = torch.triu(
            torch.ones(sequence_length, sequence_length,
                       dtype=torch.bool, device=hidden_states.device),
            diagonal=1,
        )
        scores = scores.masked_fill(causal_mask, float("-inf"))

        weights = torch.softmax(scores, dim=-1)
        return weights @ value


def main():
    batch_size, sequence_length, hidden_size = 2, 4, 8
    hidden_states = torch.randn(batch_size, sequence_length, hidden_size)
    output = CausalSelfAttention(hidden_size)(hidden_states)
    print("hidden_states shape:", hidden_states.shape)
    print("output shape:", output.shape)
    assert output.shape == hidden_states.shape


if __name__ == "__main__":
    main()
\end{lstlisting}

The explicit device argument in the mask construction keeps the mask colocated
with the input tensor when the program is moved from CPU to another device. It
does not change the lesson's mathematics.

\section{Shape analysis}
For the example configuration, the complete tensor flow is:

\begin{center}
\begin{tabular}{lll}
\toprule
Tensor & Operation & Shape \\
\midrule
$H$ & input hidden states & $2\times4\times8$ \\
$Q,K,V$ & three linear projections & $2\times4\times8$ each \\
$QK^\mathsf{T}$ & pairwise similarities & $2\times4\times4$ \\
$A$ & masked row-wise softmax & $2\times4\times4$ \\
$AV$ & weighted value mixture & $2\times4\times8$ \\
\bottomrule
\end{tabular}
\end{center}

The middle dimension of the score tensor is the querying position and the last
dimension is the position being considered. This explains why attention scores
have two sequence dimensions: every position compares itself with every other
position before the causal restriction is applied.

\section{A small tensor illustration}
The following calculation removes the batch dimension temporarily so that every
number can be seen. We use three positions and two features per position. Assume
that the projections have already produced
\[
Q=K=V=
\begin{bmatrix}
1&0\\
0&1\\
1&1
\end{bmatrix},
\qquad Q,K,V\in\mathbb{R}^{3\times2}.
\]
The rows are positions $0$, $1$, and $2$; the columns are feature coordinates.
This is a deliberately small example, not a claim that a trained model produces
these exact values.

\subsection{1. Compare every query with every key}
The key is transposed so that its feature dimension is on the inside of the
matrix multiplication:
\[
K^\mathsf{T}=
\begin{bmatrix}
1&0&1\\
0&1&1
\end{bmatrix},
\qquad
QK^\mathsf{T}
=
\begin{bmatrix}
1&0&1\\
0&1&1\\
1&1&2
\end{bmatrix}.
\]
For example, the entry in row $1$, column $2$ is the dot product
\[
Q_1\cdot K_2=[0,1]\cdot[1,1]=0\cdot1+1\cdot1=1.
\]
The row index identifies the querying position, while the column index identifies
the key being considered. The matrix multiplication is therefore
\[
(3\times2)(2\times3)=3\times3.
\]

\subsection{2. Scale and apply the causal mask}
Because the feature size is $D=2$, divide the score matrix by $\sqrt{2}$:
\[
\widetilde L\approx
\begin{bmatrix}
0.707&0&0.707\\
0&0.707&0.707\\
0.707&0.707&1.414
\end{bmatrix}.
\]
Position $0$ may see only position $0$; position $1$ may see positions $0$ and
$1$; and position $2$ may see all three positions. Therefore the masked scores
are approximately
\[
\widehat L\approx
\begin{bmatrix}
0.707&-\infty&-\infty\\
0&0.707&-\infty\\
0.707&0.707&1.414
\end{bmatrix}.
\]
The $-\infty$ entries are not values that will be retrieved. They only ensure
that their softmax probabilities become zero.

\subsection{3. Turn scores into attention weights}
Applying softmax separately to each row gives approximately
\[
A=\operatorname{softmax}(\widehat L)\approx
\begin{bmatrix}
1.000&0&0\\
0.330&0.670&0\\
0.248&0.248&0.503
\end{bmatrix}.
\]
Each row sums to $1$. Notice the triangular pattern: the zeros above the
diagonal are exactly the forbidden future positions.

\subsection{4. Mix the value vectors}
Finally, multiply the weights by the values:
\[
O=AV\approx
\begin{bmatrix}
1.000&0&0\\
0.330&0.670&0\\
0.248&0.248&0.503
\end{bmatrix}
\begin{bmatrix}
1&0\\
0&1\\
1&1
\end{bmatrix}
=
\begin{bmatrix}
1.000&0.000\\
0.330&0.670\\
0.751&0.751
\end{bmatrix}.
\]
For example, the output at position $2$ is the weighted sum
\[
o_2=0.248[1,0]+0.248[0,1]+0.503[1,1]
\approx[0.751,0.751].
\]
Thus, attention does not copy one value vector in general. It uses each row of
$A$ as a recipe for mixing the value rows. The complete calculation has the
shape flow
\[
\underbrace{(3\times2)}_{Q}
\underbrace{(2\times3)}_{K^\mathsf{T}}
\to\underbrace{(3\times3)}_{\text{scores and weights}}
\underbrace{(3\times2)}_{V}
\to\underbrace{(3\times2)}_{O}.
\]

\section{Causality as an invariance property}
The mask can be tested without inspecting individual floating-point values.
Let $H$ and $H'$ agree at position $0$ but differ at positions $1,\ldots,S-1$.
Because row $0$ of the mask permits only key $0$, we have
\[
O(H)_{0}=O(H')_{0},
\]
up to floating-point precision. A later output, such as $O(H)_2$, is allowed
to change because position $2$ can attend to positions $0$, $1$, and $2$.

An executable version of this argument is:

\begin{lstlisting}
torch.manual_seed(0)
attention = CausalSelfAttention(hidden_size=4)
first = torch.randn(1, 4, 4)
second = first.clone()
second[:, 1:, :] = torch.randn(1, 3, 4)

out_first = attention(first)
out_second = attention(second)

assert torch.allclose(out_first[:, 0, :], out_second[:, 0, :])
assert not torch.allclose(out_first[:, 2, :], out_second[:, 2, :])
\end{lstlisting}

The second assertion is probabilistic in the practical sense: with randomly
initialized parameters and different later inputs, a later output is expected
to change. The first assertion is the structural guarantee supplied by the
causal mask.

\section{Discussion}
This module is a single-head, single-layer educational abstraction. Production
language models add residual connections, normalization, multiple attention
heads or latent attention variants, positional information, feed-forward
networks, and many efficiency improvements. DeepSeek-V3 additionally uses
architecture-specific attention and mixture-of-experts components. Those
features change scale and implementation details, but not the basic logic
demonstrated here: learned projections produce compatibility scores, masking
enforces the autoregressive dependency structure, and a weighted sum retrieves
context.

The implementation also omits an explicit positional encoding. Without some
source of positional information, attention is permutation-equivariant: it can
compare content but does not inherently know whether a token came first or
last. Real Transformer systems therefore include positional information in
their attention computation or hidden states.

\section{Conclusion}
The lesson establishes causal self-attention as a sequence-to-sequence tensor
operation. Starting with $H\in\mathbb{R}^{B\times S\times D}$, the module
constructs $Q$, $K$, and $V$, forms scaled pairwise scores, removes access to
future positions, normalizes the remaining scores, and returns $AV$ with the
original hidden-state shape. The masking invariance test turns the central
autoregressive requirement into an executable claim. This small construction
provides the mathematical foundation needed to study larger Transformer and
DeepSeek-style architectures.

\begin{thebibliography}{9}
\bibitem{vaswani2017}
Ashish Vaswani et al., ``Attention Is All You Need,''
\emph{Advances in Neural Information Processing Systems}, 2017.
\url{https://arxiv.org/abs/1706.03762}

\bibitem{deepseek}
DeepSeek-AI, ``DeepSeek-V3,'' source repository, 2024.
\url{https://github.com/deepseek-ai/DeepSeek-V3}

\bibitem{pytorch}
PyTorch Contributors, ``PyTorch Documentation: Multi-dimensional arrays and
neural network modules,'' 2026.
\url{https://pytorch.org/docs/stable/}
\end{thebibliography}

\end{document}
